\section{How much nonabelian information survives a support budget?}

Let $L/\mathbb Q$ be a finite Galois extension with group $G$ of order $g$, and let $N\trianglelefteq G$ have order $n$. We regard the quotient $G/N$ as already understood. For $N=[G,G]$ the fixed field $L^N$ is abelian over $\mathbb Q$, so Dirichlet GRH gives RH for its Dedekind zeta function. To apply the one-sided method one would like an observable depending on few Frobenius types. The size of its support below counts group elements, so the corresponding density of unramified rational primes is the size divided by $g$.

The incoming packet \texttt{grh\_sparse\_singular\_moduli.zip}, in \texttt{frobenius.tex}, establishes the optimal support and margin for detecting \emph{all} new characters, as well as an identity for the zero set. The question here allows the observable to discard some characters. The answer is a finite optimization over normal subgroups and determines which intermediate fields can still occur.

Call an irreducible complex character $\chi$ new if its representation is nontrivial on $N$. Consider real, inversion-invariant class functions $f$ satisfying $f(1)=0$, $0\le f\le1$ and $c_\chi=\langle f,\chi\rangle\le0$ for every new $\chi$. The coefficients are real by inversion invariance. Define the mass of the active characters by
\[
 W(f)=\sum_{\chi\ {\rm new},\ c_\chi<0}\chi(1)^2.
\]
The summands are the dimensions of the corresponding parts of the regular representation. No analytic hypothesis is needed for this optimization.

\medskip\noindent\textbf{Support-budget theorem.}
For every $0\le B\le g$,
\[
 \max_{|\operatorname{supp} f|\le B}W(f)
 =\max_{\substack{K\trianglelefteq G,\ K\subseteq N\\ n-|K|\le B}}
      \left(\frac{g}{|K|}-\frac{g}{n}\right).
\]
For each feasible $K$, the indicator $f_K=\mathbf1_N-\mathbf1_K$ attains the value attached to $K$ using exactly $n-|K|$ elements. Every admissible $f$ factors through the quotient by its zero set in $N$.

\medskip\noindent\emph{Proof.}
For $h\in N$ all quotient characters have value $\chi(1)$. By subtracting the expansion of $f(1)$ from that of $f(h)$ we get
\[
 f(h)=\sum_{\chi\ {\rm new}}(-c_\chi)(\chi(1)-\operatorname{Re}\chi(h)).
\]
Each summand is nonnegative by unitarity. The bracket vanishes exactly when $h$ acts as the identity on the representation. Thus
$K=\{h\in N:f(h)=0\}$ is the intersection of $N$ with the kernels of all active characters. It is normal in $G$. Every character appearing in $f$ kills $K$, since this holds both for the active characters and for those inflated from $G/N$. Consequently $f$ factors through $G/K$. Its support on $N$ has size $n-|K|$, and all its active characters are characters of $G/K$ outside $G/N$. The sum of their squared degrees is at most $g/|K|-g/n$ by the regular representation degree formula. Conversely, averaging a representation over a normal subgroup projects to its invariant subspace, which is either the entire representation or zero. It follows that the new-character coefficients of $f_K$ are $-|K|\chi(1)/g$ if $K$ acts trivially and zero otherwise. This proves attainment and the theorem.\hfill$\square$

\paragraph{Two exact obstructions.} For $G=S_4$, $N=A_4$, the normal subgroups of $G$ contained in $N$ are $A_4$, the Klein four-group $V_4$ and $1$. This list also follows by inspecting the possible unions of conjugacy classes of sizes $1,3,8$ and applying Lagrange's theorem. The maximum $W(f)$ is $0$ for $B<8$, $4$ for $8\le B<11$, and $22$ for $B\ge11$. At the middle value only the two-dimensional representation inflated from $S_4/V_4\simeq S_3$ can occur. The indicator of the eight $3$-cycles attains the bound. Adding the three double transpositions attains the upper step. For any nonabelian simple $G$ and $N=G$, one active character already forces support on all $g-1$ nonidentity elements. For $A_5$ this gives an all-or-nothing threshold of $59$, even if the aim is to see a single irreducible. These obstructions use the stated sign condition on every new coefficient. Arithmetic relations between specific $L$-functions can give other criteria.

\subsection*{The cost of allowing both Fourier signs}
A sparse observable might trade strictly one-sided Fourier response for a small contribution of the opposite sign. On a nonabelian simple group the loss incurred by omitting even one additional class can be bounded sharply. For $\chi\ne1$ put
\[
 \ell_\chi(h)=1-\operatorname{Re}\chi(h)/\chi(1),\qquad
 P_\pm(f)=\sum_{\chi\ne1,\ \pm c_\chi>0}|c_\chi|\chi(1).
\]
Since each nontrivial irreducible is faithful, all losses $\ell_\chi(h)$ for $h\ne1$ are positive. Define
\[
 \kappa(G)=\min_{h\ne1}
 \frac{\min_{\chi\ne1}\ell_\chi(h)}{\max_{\chi\ne1}\ell_\chi(h)}.
\]
If a nonzero real, inversion-invariant observable $0\le f\le1$ vanishes at $1$ and at some $h\ne1$, then $P_+(f)\ge\kappa(G)P_-(f)$. Indeed $f(h)-f(1)=0$ equates the negative and positive Fourier masses weighted by their respective losses at $h$. For $\kappa(G)<1$ the bound is sharp. Choose $h,\chi,\psi$ realizing the minimum ratio, and take a positive multiple of $f(u)=\ell_\chi(u)-\kappa(G)\ell_\psi(u)$. This function is nonnegative everywhere by the definition of $\kappa(G)$, and is zero at $h$ and the identity. Its two character pairs (with complex conjugates included as needed) are disjoint because their real normalized traces differ at $h$; hence $P_+/P_-=\kappa(G)$. Scale $f$ to have maximum one.

For $A_5$ the usual character table gives the minimum-to-maximum loss ratios $3/5$, $5/8$, $(10-2\sqrt5)/15$ for elements of orders $2,3,5$, respectively. At one class of $5$-cycles the extreme normalized traces are $(1+\sqrt5)/6$ for a representation of degree three and $-1/4$ for the degree-four representation. Thus omitting even one nonidentity class costs at least $(10-2\sqrt5)/15$ in the ratio $P_+/P_-$.

\paragraph{The complete rational $A_5$ calculation.}
If $f$ must be determined by the orders of Frobenius elements, write $a,b,c$ for its nonnegative values at orders $2,3,5$. These are the three possible nonidentity cycle types in the natural degree-five action. The two conjugacy classes of order five are pooled. Their character columns are exchanged by $\sqrt5\mapsto-\sqrt5$. The \href{https://www.math.rwth-aachen.de/~Thomas.Breuer/ctbllib/ctbltoc/data/A5.html}{ATLAS character table in GAP} gives the character rows on the classes $(1,2,3,5a,5b)$, with class sizes $(1,15,20,12,12)$, as follows
\[
 \begin{array}{c|rrrrr}
 &1&2&3&5a&5b\\\hline
 \chi_1&1&1&1&1&1\\
 \chi_3&3&-1&0&\phi&1-\phi\\
 \chi_{3'}&3&-1&0&1-\phi&\phi\\
 \chi_4&4&0&1&-1&-1\\
 \chi_5&5&1&-1&0&0
\end{array}
\]
where $\phi=(1+\sqrt5)/2$. Combining the two equal degree-three coefficients, the signed contributions $60\chi(1)c_\chi$ to the masses are
\[
 v_3=-90a+72c,\qquad v_4=80b-96c,\qquad v_5=75a-100b .
\]
The first term combines both degree-three representations. Sum the positive parts of these three numbers to obtain $60P_+$ and their negative parts to obtain $60P_-$. Every nonzero admissible weight has $P_->0$ since $\sum_i v_i=-15a-20b-24c$. The full frontier for the minimum ratio under a budget of $B$ group elements is
\[
\begin{array}{c|c|c}
 B & \min P_+/P_- & (a,b,c)\ \text{attaining it}\\ \hline
 0\text{--}14 & \text{no nonzero weight} & -\\
 15\text{--}19 & 5/6 &(1,0,0)\\
 20\text{--}23 & 4/5 &(0,1,0)\\
 24\text{--}34 & 3/4 &(0,0,1)\\
 35\text{--}38 & 2/3 &(1,3/4,0)\\
 39\text{--}43 & 5/8 &(4/5,0,1)\\
 44\text{--}58 & 3/5 &(0,1,5/6)\\
 59\text{--}60 & 0 &(1,1,1)
\end{array}
\]
To verify exhaustiveness, supporting only one order fixes its ratio. Below budget $59$ at most two orders can be supported. Normalize their values to sum to one. On each interval where the signs of the three $v_i$ are fixed, their ratio of linear forms is monotone or constant. Its minimum lies at an axis or a zero of some $v_i$. The three interior candidate rays are $(4,3,0)$, $(4,0,5)$ and $(0,6,5)$, yielding the table by substitution. At budget $59$ the full indicator gives $P_+=0$. The included script recomputes the rational frontier by exact arithmetic. In particular, monitoring fewer than all three nonidentity polynomial splitting types forces a positive contribution amounting to at least three fifths of the negative mass. This statement concerns the spectrum of the arithmetic observable; converting mixed Fourier signs into a conclusion on zeros needs additional information about their multiplicities.

\paragraph{Which zero problem does the optimal observable see?}
Write $M=L^K$, $F=L^N$, $m=[M:\mathbb Q]$, $a=[F:\mathbb Q]$. Suppose RH for $\zeta_F$ and fix $K\ne N$. The prime event $f_K(\operatorname{Frob}_p)=1$ means that $p$ splits completely in $F$ and fails to split completely in $M$. Write $P$ for this set of primes, omitting primes ramified in $L$, and $\delta=1/a-1/m$. Define
\[
 R_K(x)=\sum_{p\in P,\,p\le x}\log p\,(1/p-1/x)-\delta\log x-c,
\qquad B_K(s)=\sum_{p\in P}(\log p)p^{-s}.
\]
Here $c$ is the constant term of $B_K$ at $s=1$, less $\delta$. There exists $C$ with $R_K(x)\ge -C x^{-1/2}$ for all sufficiently large $x$ if and only if $\zeta_M$ satisfies RH. The ordinary Frobenius test has no access to data distinguishing individual lifts in $G$ of an element of $G/K$.

\medskip\noindent\emph{Proof.}
The regular character of a quotient counts its identity Frobenius class. Euler factors therefore give, on $\operatorname{Re}s>1$,
\[
 B_K(s)=-\frac1a\frac{\zeta_F'}{\zeta_F}(s)
         +\frac1m\frac{\zeta_M'}{\zeta_M}(s)+H(s).
\]
Here $H$ is holomorphic on $\operatorname{Re}s>1/2$, since deleting powers $p^j$, $j\ge2$, is absolutely convergent there and only finitely many ramified Euler factors are removed. The Mellin integral $\int_1^\infty R_K(x)x^{-s}\,dx$ equals
\[
 \frac{B_K(s)}{s(s-1)}-\frac{\delta}{(s-1)^2}-\frac{c}{s-1}.
\]
The pole at one cancels. At a zero $\rho$ of $\zeta_M$ with $1/2<\operatorname{Re}\rho<1$ this expression has residue $\operatorname{ord}_\rho\zeta_M/[m\rho(\rho-1)]$. In particular every real pole has negative residue. If the proposed lower bound holds, the transform of $R_K(x)+Cx^{-1/2}$, with a bounded initial interval omitted, integrates a nonnegative function. If its convergence abscissa were greater than $1/2$, Landau's theorem would force a real singularity there. The integral is positive on the real axis to the right of the abscissa, ruling out a negative-residue pole, and the continuation is regular on $[1,\infty)$. Consequently this transform is holomorphic on $\operatorname{Re}s>1/2$. A zero of $\zeta_M$ in that half-plane would produce a pole, so the functional equation proves RH.

For the forward estimate, under RH for each of $E=F,M$, the smoothed prime ideal theorem of Greni\'e and Molteni gives
$U_E(y)=\int_1^y(\psi_E(t)-t)\,dt=O_E(y^{3/2})$.
Another integration by parts gives $\int_x^\infty (\psi_E(t)-t)t^{-2}\,dt=O_E(x^{-1/2})$. Hence the prime-power ramp for $-\zeta_E'/\zeta_E$ equals $\log x$ plus a constant and $O_E(x^{-1/2})$. Omitting ramified primes changes the constant with residual $O_L((\log x)/x)$. For the higher-power terms use $\sum_{p^j\le y,\,j\ge2}\log p\ll\sqrt y$ from the elementary Chebyshev bound, first for squares, then for the $O(\log y)$ exponents at least three. Integration against $t^{-2}$ shows that deleting them changes the constant with residual $O(x^{-1/2})$. Finally the Mellin identity identifies the constant as $c$. This proves the two-sided endpoint bound.\hfill$\square$

The analytic reasoning builds on Suzuki's one-sided prime-sum criterion and the earlier incoming packets on relative races and prime sampling. Its role here is to identify the field left visible at each step of the support theorem. Establishing any required new lower bound for primes still needs arithmetic input. See \href{https://arxiv.org/abs/2411.07436}{M. Suzuki, \emph{On variants of Chebyshev's conjecture}} and \href{https://doi.org/10.1090/mcom3039}{L. Greni\'e and G. Molteni, \emph{Explicit smoothed prime ideals theorems under GRH}}.
